\section{Teil 4}
\begin{block}[Paradies]
        Vom Paradies bis zum Sündenfall und der Vertreibung aus dem Paradies haben wir uns an den letzten Sonntagen beschäftigt. Da haben wir gelernt, dass Gott uns Menschen einen Auftrag gegeben hat.
        \begin{bibelbox}{SCHL}{IMos}{1:28}
            Und Gott segnete sie; und Gott sprach zu ihnen: Seid fruchtbar und mehret euch und füllt die Erde und macht sie euch untertan; und herrscht über die Fische im Meer und über die Vögel des Himmels und über alles Lebendige, das sich regt auf der Erde.
        \end{bibelbox}
        Das ist ein Auftrag an uns. Das Gebot, welches Gott den beiden Menschen damals gab, nicht von der Frucht vom Baum der Erkentniss von Gut und Böse zu essen, wurde nicht eingehalten. Zur Strafe mussten sie das Paradies verlassen, es kam der Tod und die Sünde ins Leben. Der Satan wurde der König der Welt. Den Auftrag, die Erde zu bevölkern und zu bebauen, gilt immer noch, aber durch den Sündenfall wurde die Ausführung erschwert. Wir müssen jetzt im Schweisse des Angesichts die Erde bebauen und die Frauen müssen unter Schmerzen gebären.
        
        Der Satan ist nun der Fürst dieser Welt. Und Gott sagt uns in Vers 15 wie er diesen Fürsten bekämpfen will.
        \begin{bibelbox}{SCHL}{IMos}{3:15}
            Und ich will Feindschaft setzten zwischen dir und der Frau; zwischen deinem Samen und ihren Samen: Er wird dir den Kopf zertreten und du wirst ihn in die Ferse stechen.
        \end{bibelbox}
        Jesus sagt dann in 
        \begin{bibelbox}{SCHL}{Joh}{12:31-32}
            Jetzt ergeht ein GEricht über die Welt. Nun wird der Fürst dieser Welt hinausgeworfen werden; uns ich, wenn ich von der Erde erhöht bin, werde alle zu mir ziehen.
        \end{bibelbox}
        Heute wissen wir, dass dieser Fürst unser Herr Jesus Christus ist. Damals wussten sie zwar, dass jemand kommen wird, aber sie wussten nicht, wer er ist oder wann er kommt.
        Hier kommen wir jetzt zu der Phase 3 in Gottes Plan sein Königreich aufzurichten. Diese Phase dauert in der Bibel von der Vertreibung aus dem Paradies bis zu Maleachi und zum Auftritt von Jesus Christus. Eine Dauer von ca. 4000 Jahren laut der Zeitrechnung von Roger Liebi.

        In dieser Zeit hat Gott aber nicht tatenlos zugesehen wie die Menschheit verkommt und ins Verderben rennt, sondern er hat einen Plan, den er durchzieht. Und dieser Plan wird uns in den nächsten Sonntagen beschäftigen. Es gibt verschiedene Bündnisse, die Gott, mit uns Menschen geschlossen hat, um zum Ziel zu kommen. Das Kommen von Jesus unserem Erlöser wird im Laufe der Zeit immer konkreter.
        \begin{itemize}
            \item In \bibleverse{IMos}(3:15) wurde Adam und Eva gesagt, dass irgendwann einer kommt, welcher der Schlange den Kopf zertreten wird und von einer Frau abstammen wird. Da sie die einzige Frau war, dachte sie wahrscheinlich, dass ihr erster Sohn das sein wird.
            \item Nach der Sintflut lässt der Vers \bibleverse{IMos}(9:26) erahnen, dass die Linie des Retters über Sem läuft.
            \begin{bibelbox}{SCHL}{IMos}{9:26}
                Gepriesen sei der Herr, der Gott Sems, und Kaanan sei sein Knecht.
            \end{bibelbox}
            \item \bibleverse{IISam}(7:16) sagt uns, dass die Linie ganz sicher über die Stammlinie von König David gehen wird.
            \begin{bibelbox}{SCHL}{IISam}{7:16}
                dein Haus und dein Königreich sollen ewig Bestand haben wie dein Angesicht; dein Thron soll auf ewig fest stehen.
            \end{bibelbox}
            Beim Lesen dieses Bibelverses ist mir so richtig bewusst geworden, dass Gott \betonung{ewig} sagt. Und ewig ist ewig, nicht nur bis Jesus wiederkommt. 
        \end{itemize}
        Um sein Ziel mit uns Menschen zu erreichen, hat ER auf dem Weg verschiedene Bündnisse geschlossen und sich ein eigenes Volk aufgebaut. Diese Bündnisse werden wir in den nächsten Sonntagen zusammen anschauen.        
    \end{block}
    \numcirc{1}
    \begin{block}[Noah Bund]
        Finden wir in \bibleverse{IMos}(8:21-22)(9:8-17)        
    \end{block}
    \numcirc{2}
    \begin{block}[Abraham Bund]
        \bibleverse{IMos}(12:1-3)
    \end{block}
    \numcirc{3}
    \begin{block}[Mose Bund]
        Oder auch der Bund am Sinai, Alter Bund oder Gesetzesbund genannt. Dieser Bund wird mit den Gesetzestafeln in \bibleverse{IIMos}(34:1) oder die ersten Tafeln \bibleverse{IIMos}(31:18)(32:16) ins Leben gerufen.
    \end{block}
    \numcirc{4}
    \begin{block}[David Bund]
        Mit diesem Bund versricht Gott, dass der Retter aus der Stammline der Nachkommen Davids kommen wird. \bibleverse{IISam}(7:12-16)(23:5)
    \end{block}
    \numcirc{5}
    \begin{block}[Neuer Bund]
        Von dem Lesen wir in \bibleverse{Jer}(31:27-40), \bibleverse{Lk}(20:17-22)
        Und nachher im Abendmahl gedenken wir an diesen neuen Bund.
    \end{block}
    So jetzt habt ihr Stoff zum Lesen und Studieren, bis Janina und ich wieder von Israel zurück sind.
